% HD-S31. GPT-6 Astra Pro, /root/sixfolds; inherited effort not exposed.
% 21 November 2026. Published classification input, exact quotient maps.
\subsection{HD-S31: the rank-three block endpoints and the \(G_{24}\) determinant}
\label{sec:rank-three-endpoints}

For the actual rank-three block systems in the last paragraph of
HD-S30, a sixfold period differential of rank nine has one of the
following coefficient sources:
\[
\begin{array}{c|c|c}
\text{coefficient group}&K&\text{ordinary determinant information}\\\hline
G(d,d,3),\ d=3,4,6&\mathbb Q(\zeta_d)&[-1]\\
G_{24}&\mathbb Q(\sqrt{-7})&[-1]\\
G_{25},G_{26}&\mathbb Q(\sqrt{-3})&
 \text{within the already covered Eisenstein field}.
\end{array}
\]
The last line makes no assignment of determinant class to an arbitrary
tuple. All Eisenstein sixfold determinant classes already have the
separate proved construction of HD-G10--HD-G11. Thus this table closes
the indicated rank-three block route to every remaining field and to
every remaining non-Eisenstein nonsplit determinant; it does not
classify representations without the block hypotheses of HD-S30.

\paragraph{The trace-field map and the precise published list.}
Keep the exact induced-module isomorphism (S30.1). Its character is
the induced character of W, so its trace at g is a sum of traces of
W on the fixed cosets, with zero contribution from the other cycles.
Consequently
\[
 K=\mathbb Q(\operatorname{tr}_V G)
 \ \subseteq\ \mathbb Q(\operatorname{tr}_W H)
 \ \subseteq\ K.
\]
The last inclusion follows from the actual K-subspace W. Thus the
coefficient representation has character field exactly K. The
rank-nine argument of HD-S30 gives the actual coefficient quotient
\(X\to Y=\mathbb P^1\), with twelve local reflections generating
its finite, faithful, irreducible rank-three image \(\Gamma\).

The published classification used at this point is Cohen,
\emph{Finite complex reflection groups}, Theorem (2.4), pp.386--387,
for imprimitive groups, and Theorem (5.12), p.435, together with
the rank-three rows of (4.16), p.412, for nonreal primitive groups.
The relevant fields are in Marin--Michel, \emph{Automorphisms of
complex reflection groups}, Table 1, author-version pp.20--21
(1 February 2010), with its definition of character field on p.2:
\[
\begin{array}{c|ccccc}
\Gamma&G_{23}&G_{24}&G_{25}&G_{26}&G_{27}\\\hline
\mathbb Q(\operatorname{tr}\Gamma)&\mathbb Q(\sqrt5)&
\mathbb Q(\sqrt{-7})&\mathbb Q(\zeta_3)&
\mathbb Q(\zeta_3)&\mathbb Q(\zeta_3,\sqrt5).
\end{array}
\]
The same rows were checked in Jean Michel's primary four-page group
table, pp.1 and 3. A real field \(\mathbb Q(\sqrt5)\) cannot
equal an imaginary quadratic field, and the displayed field for
\(G_{27}\) has degree four: its real quadratic subfield is
distinct from its imaginary quadratic subfield. These eliminate
those two rows by the proved field equality, leaving exactly the
three primitive possibilities in the first table.

For the monomial groups \(G(d,e,3)\), the entries are in
\(\mathbb Q(\zeta_d)\). Conversely the element
\[
 \operatorname{diag}(\zeta_d,\zeta_d^{-1},1)\,(23)
                  \in G(d,e,3)
\]
has trace \(\zeta_d\); its total diagonal color is zero, so
it belongs for every divisor e of d. Thus its character field is
exactly \(\mathbb Q(\zeta_d)\), by actual matrices.
Equality with imaginary quadratic K forces \(d=3,4,6\).

\paragraph{The monomial endpoint is an actual unramified genus-four cover.}
The conjugacy to the displayed monomial K-model can be made over K:
solve its intertwining equations over K. Their complex solution space
is one-dimensional and contains an invertible map, so any nonzero
K-solution is invertible. Its images of the standard coordinate lines
are therefore actual K-lines in W. This gives the specified line
system and character used in the following quotient maps.
Let \(H_1\subset\Gamma\) stabilize the first coordinate line
and let \(\chi:H_1\to\mu_d\) be its scalar action. It has
full order d, because the displayed diagonal element lies in
\(H_1\) and has scalar \(\zeta_d\). The quotient maps are
\[
 X\longrightarrow Z'=X/\ker\chi
       \xrightarrow{\pi}Z=X/H_1\longrightarrow\mathbb P^1,
 \qquad \deg(Z\to\mathbb P^1)=3.
\]
The cover \(\pi\) is the actual connected cyclic cover of order
d. Applying (S30.1) to its one-dimensional coordinate blocks gives
the same exact invariant-stalk Shapiro map and polarized coefficient
comparison. Thus its primitive character period is the coefficient
period being tested, with a fixed rational polarization multiplier.

Write \(t=g(Z)\) and let r be the true branch number of \(\pi\).
For a rank-one faithful character every nonidentity local monodromy
has codimension one, so the six-dimensional multiplicity gives
\[
                       6=2t-2+r,\qquad r=8-2t.
\]
The actual cyclic-cover family over the stable pointed curve
\((Z,D)\) has dimension at most
\(3t-3+r=t+5\), by the same discrete-monodromy deformation
map as in HD-S30. Here \(0\le t\le4\); the unmarked case
occurs only at \(t=4\), which is stable as well. Rank nine
therefore forces \(t=4,r=0\).

This also determines the monomial group. A diagonal reflection
causes a nontrivial scalar monodromy on its single colored fixed
sheet, whereas a transposition reflection has complementary colors
on its two-cycle and hence trivial scalar monodromy after that
cycle is traversed. Thus r=0 excludes every diagonal reflection.
All twelve generating reflections have total color zero, forcing
\(e=d\). Finally HD-S01--HD-S02 apply to every such unramified
cyclic cover: put its primitive monodromy on the first symplectic
handle; each of the other three handles lifts as a regular cyclic
module. On its primitive part, the displayed a- and b-cycle
generators A and B have Hermitian Gram matrix
\[
 \begin{pmatrix}0&\gamma_d\\-\gamma_d&0\end{pmatrix},
 \qquad\gamma_3=\sqrt{-3},\quad\gamma_4=2i,\quad
 \gamma_6=2\sqrt{-3},
\]
for the restricted theta polarization. In the sixth-order case the
actual integral generators are \(S(\sigma)a,\sigma S(\sigma)a\)
and the corresponding b-cycles, with
\(S(t)=t^4+t^3-t-1\), as in HD-S02. The explicit K-basis
change \((A,B)\mapsto(A,B/\gamma_d)\) has Gram matrix U.
In three blocks this is a rational K-isometry to \(U^3\), so
the ordinary determinant class is \([-1]\) for every cover.
The degree-27 quotient in HD-S02 is an additional integral-lattice
construction: for d=3 it descends the restricted theta form, and
for d=6 it descends half that form, via the identity on the real
vector space and an index-27 overlattice. The present determinant
argument needs only the explicit rational basis change above; it
does not identify the different cover moduli.

\paragraph{Every twelve-reflection \(G_{24}\) endpoint is split.}
Use the exact \(G_{24}\) model of HD-S13. With
\(\lambda^2+\lambda+2=0\), its coefficient lattice
\(\mathcal O^3\), \(\mathcal O=\mathbb Z[\lambda]\), has
positive Hermitian Gram matrix
\[
 H_0=\begin{pmatrix}2&-1&-\lambda\\
                   -1&2&-1\\1+\lambda&-1&2\end{pmatrix},
                  \qquad\det H_0=1.
\]
The group contains \(-I\), and all its reflections are
involutions. These facts were verified on the full 336-element
matrix closure in HD-S13, independently of a particular tuple.
Any K-defined model of the same reflection representation is
related to this one by an actual K-linear intertwiner: solve the
linear equations for the intertwiner over K. After extension to
\(\mathbb C\) their solution space has dimension one and
contains an invertible matrix, so a nonzero K-solution is invertible.
The two embeddings are interchanged by conjugation, which leaves
the real ordinary determinant class unchanged. Invariant positive
Hermitian coefficient forms differ by a positive rational scalar,
by the same one-dimensional commutant calculation. This scalar
is a norm after taking the determinant in parabolic rank six.

Put \(R=\mathcal O[1/2]\). This is a PID, by the explicit
norm-Euclidean proof for \(\mathcal O\) in HD-S14. The lattice
with Gram \(H_0\) is Hermitian self-dual over R. For any full
generating twelve-reflection tuple, the central element \(-I\)
makes \(H_0(\Sigma,L)=H_0(\Sigma,L^\vee)=0\) on the
bordered sphere: \(-2\) is a unit. Each local averaging projector
\((1+g_j)/2\) splits invariants and coinvariants over R.

Here all the integral hypotheses of the bordered-surface proof in
HD-S27 hold with this ring. The surface retracts onto a graph, so
absolute first homology is free. Poincare--Lefschetz cap products
and the universal coefficient sequence give perfect pairings of
absolute and relative first homology, with the dual coefficients.
Boundary invariants and coinvariants are free, so the image of
absolute in relative homology is primitive, with free cokernel.
Every functional on this image extends to the relative module;
the dual absolute class maps to a unique class in the dual image,
by adjointness of the two absolute-to-relative maps. Thus the
induced pairing on the parabolic image is perfect over R. The
invariant-stalk disc calculation identifies it with
\(H^1(\mathbb P^1,j_*R^3)\), preserving the cup product.
Combining with \(H_0\) gives a perfect skew-Hermitian form S
on a free rank-six R-module.

Its determinant is fixed by conjugation, since the rank is even,
and is an R-unit. Therefore
\(\det S\in R^\times\cap\mathbb Q=\{\pm2^a:a\in\mathbb Z\}\).
Set \(s=\sqrt{-7}\). The Hermitian cup form is \(sS\), and
\[
                    \det(sS)=-7^3\det S.
\]
Its signature is \((3,3)\), so this determinant is negative
and \(\det S=2^a>0\). Both 7 and 2 are norms:
\(N(s)=7\), \(N(\lambda)=2\). Consequently
\[
              \det(sS)=-N\bigl(s^3\lambda^a\bigr),
\]
including negative integers a. Its ordinary determinant class
is exactly \([-1]\). The actual coefficient/Jacobian map in
HD-S13 and the real pairing calculation (S30.3) multiply the form
by a fixed nonzero rational scalar; its sixth power is a norm.
Every K-linear change of lattice basis also multiplies the
determinant by a norm. Hence the actual restricted-theta sixfold
has the same split class, for every such tuple.

For a rank-nine family, changing a compatible polarization through
a fixed isogeny cannot create a different ordinary determinant
class at its general point. The isogeny pulls the target form back
to a rational compatible form on the same marked Hodge space.
Rank nine gives an open in its nine-dimensional unitary domain.
The explicit centralizer calculation of HD-S17 shows that a rational
endomorphism preserving all those Hodge subspaces is in K; its
self-adjoint part is \(\mathbb Q\). Thus the new form is a
positive rational multiple of the old form. Its rank-six determinant
changes by the norm of the third power of that multiple. This is
the actual polarization transport, with no unspecified comparison
between different moduli problems.

\paragraph{Proof audit and remaining scope.}
The complete author proof has been read after writing. The published
classification inputs are limited to the exact ranks, list and fields
specified above. The induced trace-field map, both quotient towers,
the period-rank reductions, and the universal \(G_{24}\) determinant
are proved here using the previously audited explicit coefficient and
lattice maps. No determinant was assigned by sampling tuples, and no
new dominance claim is made for any allowed endpoint. Independent
review is recorded separately when it occurs.
